\begingroup
\selectlanguage{english}
\noindent\textbf{Attachments, ordinary direct images, and homotopy.}

These pages, Chapter~XII, section~V, manuscript pp.~59--122, begin with a concrete problem: how to attach one category to another along a full subcategory, and how to recognize when the attachment respects weak equivalences. The first part answers the categorical question completely. The second discovers how much more is needed for the homotopical question. Along the way Grothendieck introduces total asphericity, isolates the Dwyer condition, and reaches the Grothendieck construction as a replacement for an ordinary pushout. The text is exploratory in a mathematically consequential sense. Some tentative arguments are explicitly retracted; one proposed theorem survives in the text even though its asserted reverse implication is false. The commentary keeps these different situations separate.

An ``open'' subcategory $U\subseteq X$ is a sieve: an arrow whose target lies in $U$ has its source there as well. Its complementary full subcategory $T$ is a cosieve. There are no arrows from $T$ to $U$, so the information not already contained in these two categories is exactly the bifunctor
\[
 H:U^{\mathrm{op}}\times T\longrightarrow\mathrm{Set},
 \qquad H(u,t)=\Hom_X(u,t).
\]
Its two functorial actions encode composition on either side of a crossing arrow. Conversely any such bifunctor constructs a category on the disjoint union of the object sets. This is the collage of a profunctor. Maps and natural transformations out of the collage have equally explicit descriptions. The formulas on manuscript p.~61 must be read relatively: changing $U$ and $T$ by pullback of $H$ gives a cartesian lifting over the category of pairs of categories, not the unqualified representability asserted in the printed displays. The corrected universal property is proved in the commentary.

This description also explains the relation with open--closed gluing of topoi. A presheaf on the collage is a pair of presheaves on $U$ and $T$, together with the restriction data across $H$. The gluing functor is
\[
 \varphi(F)(t)=\Hom_{\widehat U}(H(-,t),F).
\]
It preserves limits and is right adjoint to the cocontinuous extension of $t\mapsto H(-,t)$. The converse recovers the collage from an accessible gluing functor preserving small limits. This is already the construction of SGA~4, Expos\'e~IV, Exercise~9.5.10; the comparison is exact, including the characterization of the resulting topos as a presheaf topos \citep{SGA4I1972}. Natural transformations between right adjoints reverse the direction of transformations between their left adjoints, a convention that must be retained when stating the category-level equivalence. The historical connection helps explain why the categorical attachment problem first appears here in the language of an open and its closed complement.

The advantage of the collage is that a pushout can be computed by an ordinary Kan extension. Given $p:U\to U'$ and $r:T\to T'$, the new crossing bifunctor is
\[
 H'=\operatorname{Lan}_{p^{\mathrm{op}}\times r}H.
\]
Its pointwise coend description records a crossing arrow together with paths in the new categories to its two ends, modulo the relations imposed by composition. When only $U$ changes, this becomes
\[
 \Hom_{X'}(u',t)\cong\pi_0(u'\backslash U/t),
 \qquad X'=X\amalg_U U'.
\]
The category on the right remembers factorizations through $U$; the ordinary pushout retains only their connected components. This formula proves that $U'$ is again a sieve, that the complement remains $T$, and that the support of the gluing weight is unchanged. It is also the source of the later difficulty: taking components does not preserve the higher homotopy of a factorization category.

\medskip
\Needspace{5\baselineskip}
\noindent\textbf{The gluing theorem that survives, and the reverse implication that fails.}

If a category is the union of two sieves, every simplex of its nerve lies in one of them. Its nerve is therefore the simplicial union along the intersection, and the square is a homotopy pushout. Consequently weak equivalences on the intersection and on both pieces induce a weak equivalence on the union. This is the valid gluing direction. It does not imply the symmetric rule W(7), which asks that any two among equivalence on the intersection, equivalence on both pieces, and equivalence on the union imply the third.

Theorem~2 on manuscript p.~78 asserts that this stronger rule holds for $W_\infty$. A complete counterexample uses a connected acyclic noncontractible finite complex. Take a finite poset $A$ with that homotopy type and adjoin two incomparable terminal cone points. The two one-cone subcategories are sieves, each contractible, and their union has the homotopy type of the suspension of $BA$. The suspension is simply connected and acyclic, hence contractible, whereas $BA$ remains noncontractible. Thus the map from a cone to the union is a weak equivalence, but the opposite map from the intersection to the other cone is not. The counterexample refutes both the reverse part of W(7 bis) and the full W(7); it does not merely exhibit a missing proof.

The commentary gives the classical two-cell input explicitly, verifies its homology and a nontrivial fundamental-group quotient, and passes to a finite simplicial model and its face poset \citep[Example~2.38 and Theorem~2C.5]{Hatcher2002}. A further development supplies finite tests against every bounded homotopy-group version of the axiom. It also explains why the proposed Mayer--Vietoris argument succeeds for fixed-coefficient homology equivalences: the three mapping-cone complexes form a short exact sequence. The failure occurs in passing from that homological statement to the asserted unstable homotopical converse.

This finding changes how the subsequent pages must be read. Assertions deduced under W(7) or W(7 bis) remain conditional assertions. Their application to $W_\infty$ cannot use Theorem~2. In particular, the proposed one-point converse to asphericity on pp.~84--85 is not an unconditional theorem about weak homotopy equivalences. The audit records the actual hypotheses and, where possible, supplies a direct proof of a useful consequence that avoids the false reverse principle.

\medskip
\Needspace{5\baselineskip}
\noindent\textbf{Total asphericity measures ordinary images.}

For a presheaf $F$ on $U$, total $W$-asphericity means that the category of elements of $p_!F$ is $W$-aspherical for every functor $p:U\to U'$. These are ordinary set-valued Kan extensions. For a category $V$, the definition applies to its terminal presheaf. It is substantially stronger than $V\to e\in W$, and it is not interchangeable with other uses of ``totally aspherical'' in the literature on test categories.

The relevant ordinary factorization sends $f:V\to V'$ to
\[
 V\longrightarrow V'/f_!e_V\longrightarrow V'.
\]
The second map is a discrete fibration; the first is final in the ordinary categorical sense. Its undercategories are nonempty and connected, rather than necessarily contractible. This is the comprehensive factorization. The commentary proves its universal property and relates it to the treatment of Street--Walters in Kelly's account \citep[\S4.7]{Kelly1982}. It yields a useful formulation developed here: a category is totally $W$-aspherical exactly when every target of an ordinary final functor out of it is $W$-aspherical. In particular the class of totally aspherical categories is closed under ordinary final images. This reformulation identifies the quantified condition without confusing ordinary finality with homotopy finality.

Grothendieck's attempted stability under unions fails because a left Kan extension need not preserve the monomorphisms describing a cover. His attempted use of retractions does not repair that defect: a categorical retraction does not supply a splitting of the associated inclusion of presheaves. A finite example makes both issues explicit. A final functor from the contractible outgoing three-object fork lands in a four-object category whose classifying space is $S^2$. Its relevant comma category is a square, connected but not contractible. Two distinct representatives of a Kan-extension value become equal after adjoining one branch. The construction was located in Grothendieck's Appendix~XII and is credited there. It also produces the pushout failures on pp.~104--107; adjoining a further cone gives a classifying space $S^3$, so their validity can be checked independently of the disputed one-point converse.

The converse from totality of a presheaf to totality of its category of elements is established here for subterminal presheaves, as the source proposes. A general answer was not located in the primary literature checked. The same limitation applies to certain composition and classification questions posed later in the excerpt. The report describes those searches; it does not turn an unsuccessful search into a claim of present-day openness.

\medskip
\Needspace{5\baselineskip}
\noindent\textbf{Good attachments and the passage to integration.}

A $W$-Dwyer sieve asks that every nonempty crossing weight be totally $W$-aspherical. Ordinary Kan-extension transitivity proves that this condition is preserved by pushout. The stronger version, expressed through totality of the categories of elements themselves, uses the final first map of the comprehensive factorization. Classical Dwyer immersions are the representable case: the inclusion of $U$ into the cosieve it generates admits a right adjoint. Filtered examples show why total asphericity is broader than the existence of a terminal object. The interval-product calculation on p.~114, however, shows directly why the classical adjoint condition is too rigid for every intended homotopy attachment.

The manuscript's $W$-cofibrations have since been studied as $h$-cofibrations, defined by preservation of weak equivalences under cobase change on undercategories. In left proper model categories they characterize attachments whose pushouts compute homotopy pushouts; the theory also recovers the equivalence between trivial $h$-cofibrations and couniversal weak equivalences \citep[Proposition~1.6, Corollary~1.9 and Lemma~1.10]{BataninBerger2017}. This is a precise later form of the useful gluing principle. For classical Dwyer maps, the nerve comparison is now known to be a weak categorical equivalence, with a generalization to discretely flat cofibrations \citep[Theorems~1.5--1.6]{HackneyOzornovaRiehlRovelli2024}. These later notions are compared with the definitions here without identifying their different ambient weak equivalences.

The final properness criterion is independent of the failed cube axiom. Under the necessary connectivity hypothesis, the inclusion of a terminal fibre is ordinary final. Its terminal presheaf therefore has direct image equal to the terminal presheaf of the total category. Composing Kan extensions proves that total asphericity of that fibre implies total asphericity of the total category. The commentary gives the general component-colimit formula and writes out the natural transformation requested in the footnote on p.~119, including its action on representatives and its compatibility with composition.

The excerpt then begins a new construction. Integrating a span retains its middle category and the arrows connecting it to both sides. For two points mapping to a point on each side, the ordinary pushout is a point, whereas the integrated category has classifying space a circle. Two categorical mapping cylinders give a direct proof that the integrated span models the homotopy pushout. The final displayed identity $\pi_!\xi\cong\mathrm{hot}_W(\mathscr X)$ states the intended derivator interpretation. The supplied pages end at that identity; the announced reduction to closed immersions continues beyond them.

\medskip
\noindent\textit{Commentary, mathematical audit, and the explicitly marked developments were produced by GPT-6 Astra (OpenAI, ChatGPT Pro mode). Source references use the manuscript pagination; the report gives its concordance with the attached PDF and the edition's printed pages.}
\endgroup
